% =============================================================================
% Seção 1 — Introdução
% =============================================================================

\section{Introdução}
\label{sec:introducao}

A esclerose lateral amiotrófica (ELA), internacionalmente denominada motor neurone disease
(MND) ou amyotrophic lateral sclerosis (ALS), é uma condição neurodegenerativa progressiva que destrói os neurônios motores
superiores e inferiores, resultando em perda de movimento voluntário, fala e, por fim,
função respiratória. A incidência global é de aproximadamente 2--3 casos por
100.000 pessoas-ano, afetando cerca de 41.000 indivíduos na Austrália e 30.000 no
Brasil~\cite{mehta2018mnd}. A sobrevida mediana após o diagnóstico é de 2--5 anos; a
insuficiência respiratória, tipicamente tratada com ventilação não invasiva (VNI), é
a principal causa de morte. Notavelmente, a VNI demonstrou prolongar a sobrevida por
um fator de aproximadamente três em comparação ao melhor tratamento farmacológico
disponível~\cite{berlowitz2021niv}, tornando o período de uso da VNI clinicamente
significativo e prolongado. Durante essa fase, os pacientes permanecem cognitivamente
intactos enquanto se tornam progressivamente incapazes de falar ou gesticular de
formas que as tecnologias de comunicação padrão consigam detectar.

Trabalhos recentes demonstraram que a assincronia paciente-ventilador (PVA) — uma
incompatibilidade entre os ciclos do ventilador e o esforço respiratório do paciente
— é um contribuinte significativo para a não adesão à VNI na ELA, causando
desconforto, fragmentação do sono e potencialmente desfechos clínicos
adversos~\cite{atkeson2011pva, yang2026pva, berlowitz2025tla}. Pacientes que
experimentam desconforto relacionado à PVA durante o uso noturno da VNI não
conseguem sinalizar angústia facilmente aos cuidadores usando as tecnologias de CAA
existentes, que requerem conectividade de rede ou hardware dedicado inadequado para
o ambiente domiciliar. O NeuroAssist AI aborda diretamente essa lacuna: ao fornecer
um canal de sinalização local e sempre disponível que requer apenas uma webcam comum,
permite que os pacientes alertem os cuidadores durante períodos sem supervisão —
incluindo enquanto utilizam a VNI à noite.

Em estágios avançados, uma proporção substancial de pacientes experimenta uma condição
semelhante ao locked-in, na qual o controle voluntário se limita a movimentos faciais
residuais sutis: pequenas contrações palpebrais, elevações de sobrancelha ou
microexpressões que persistem após a perda da função motora dos membros e orofacial.
Esses gestos residuais representam o último canal voluntário disponível para
comunicação. Sem um meio confiável de detectá-los e agir sobre eles, os pacientes não
conseguem sinalizar angústia, desconforto ou necessidades aos cuidadores — com
consequências diretas para a qualidade de vida e a segurança do paciente durante
períodos sem assistência.

Os sistemas de comunicação aumentativa e alternativa (CAA) existentes não atendem
adequadamente essa população. Dispositivos de rastreamento ocular de ponta (e.g.,
Tobii Dynavox) e softwares de CAA dedicados (e.g., Grid~3) oferecem alta precisão,
mas requerem conectividade sem fio estável, hardware dedicado custando mais de
USD~3.000, e configurações calibradas em fábrica que não conseguem se adaptar ao
perfil motor em declínio de um paciente individual~\cite{beukelman2013aac}.
Interfaces cérebro-computador (BCI) via EEG oferecem uma alternativa, mas envolvem
hardware invasivo ou volumoso inadequado para implantação domiciliar~\cite{zander2011passive}.
Abordagens baseadas em câmera de baixo custo existem na literatura~\cite{soukupova2016},
mas dependem de comparações simples de limiar por quadro que geram taxas de falsos
positivos inaceitavelmente altas quando aplicadas à atividade facial involuntária
presente em pacientes com ELA durante a VNI.

Apresentamos o \textbf{NeuroAssist AI}, uma plataforma de comunicação assistiva
offline-first e preservadora de privacidade projetada para implantação domiciliar com
uma webcam comum. O sistema transforma gestos faciais residuais em solicitações de
emergência estruturadas usando um pipeline de visão computacional WebAssembly
executado localmente, uma máquina de estados de gesto deliberado e um perfil de
detecção calibrável por paciente — sem transmitir nenhum vídeo para fora do
dispositivo.

Este artigo apresenta três contribuições originais:
\begin{enumerate}
  \item \textbf{Um pipeline de detecção offline-first e preservador de privacidade.}
    O modelo MediaPipe FaceLandmarker e seu runtime WebAssembly são servidos a partir
    de ativos locais; nenhum quadro de vídeo ou dado biométrico sai do dispositivo
    do paciente. O delegate de CPU garante compatibilidade com hardware comum,
    independentemente da disponibilidade de GPU.

  \item \textbf{Uma máquina de estados de gesto deliberado de seis estados com
    calibração específica por paciente.}
    Em vez de um único limiar EAR, a máquina de estados exige duração mínima de
    fechamento ($\geq$3 quadros, ${\approx}$100\,ms) em cada uma de duas piscadas
    sequenciais, reabertura confirmada entre elas, e uma janela de tempo configurável.
    Os limiares são armazenados em um \texttt{DetectionProfile} versionado e adaptados
    por paciente. Esse design aborda especificamente as fontes de falsos positivos
    relevantes para a ELA: piscadas involuntárias, artefatos de ptose parcial e
    deriva lenta de limiar.

  \item \textbf{Uma arquitetura extensível de microsserviços hexagonais com tipos de
    gesto plug-in.}
    O campo \texttt{engineFeatureRef} no perfil de detecção desacopla o motor de
    reconhecimento de gesto do catálogo de comandos, permitindo que novos gestos
    (elevação de sobrancelha, piscada de um olho) sejam registrados sem modificar o
    código de domínio existente. Regras ArchUnit aplicadas em tempo de compilação
    garantem que essa separação persista à medida que o sistema evolui.
\end{enumerate}

O restante do artigo está organizado da seguinte forma. A Seção~\ref{sec:relacionados}
revisa os trabalhos relacionados em CAA, detecção de gesto facial e computação de
borda para saúde. A Seção~\ref{sec:arquitetura} descreve a arquitetura do sistema em
detalhe. A Seção~\ref{sec:avaliacao} relata a avaliação experimental. A
Seção~\ref{sec:discussao} discute limitações e trabalhos futuros, e a
Seção~\ref{sec:conclusao} conclui o artigo.
